\documentclass[11pt]{article}


\setlength{\parindent}{0pt}
\usepackage{xltxtra}
\usepackage{hyperref}
\hypersetup{hidelinks}
\usepackage{url}
\urlstyle{tt}
\usepackage{xcolor}
\definecolor{CVBlue}{RGB}{23,110,191}
\usepackage{calc}
\usepackage{graphicx}
\usepackage{tikz}
\usetikzlibrary{calc}
\usepackage{fontspec}
\usepackage{xeCJK}
\usepackage{enumitem}
\CJKsetecglue{} %% 取消中文与数字之间的间隙


%% 主文档字体设置
\setmainfont[
    Path = fonts/Main/,
    Extension = .otf,
    BoldFont = texgyretermes-bold.otf, % 加粗字体
]{texgyretermes-regular.otf} % 正文字体

% 中文字体设置
\setCJKmainfont[
    Path = fonts/hansans/,
    Extension = .ttf,
    BoldFont = NotoSansSC-Bold.ttf, % 加粗字体
]{NotoSansSC-Regular.otf} % 正文字体


\usepackage{relsize}
\usepackage{xspace}

% 使用 fontawesome（部分图标）
\usepackage{fontawesome} 

% A4纸，上下左右边距
\usepackage[
    a4paper,
    left=1.2cm,
    right=1.2cm,
    top=1.5cm,
    bottom=1cm,
    nohead
]{geometry}

\renewcommand{\baselinestretch}{1.5} % 行间距设为1.5

\usepackage{titlesec}
\usepackage{enumitem}
\setlist{noitemsep} % 取消列表项间的额外间距
%\setlist{nosep} % 取消所有垂直间距
\setlist[itemize]{topsep=0.25em, leftmargin=*}
\setlist[enumerate]{topsep=0.25em, leftmargin=*}

% --- 用于控制【不同项目之间】的垂直距离 ---
\newlength{\interProjectSpacing}
\setlength{\interProjectSpacing}{0.9em} % <--- 在此调整项目之间的距离
\newcommand{\projectsep}{\vspace{\interProjectSpacing}}

% --- 用于控制【项目标题】与下方【项目描述】的距离 ---
\newlength{\intraProjectTitleSep}
\setlength{\intraProjectTitleSep}{0.4em} % <--- 在此调整标题和描述的距离
\newcommand{\titlebreak}{\\[\intraProjectTitleSep]}

% --- 用于控制【项目描述】与下方【要点列表】的距离 ---
\newlength{\intraProjectListTopSep}
\setlength{\intraProjectListTopSep}{0.2em} % <--- 在此调整描述和列表的距离

% =======================================================================


\titleformat{\section}         % 定制 \section 命令 
{\large\bfseries\raggedright} % 将 section 标题设置为大号、粗体且左对齐
{}{0em}                      % 可用于为所有 section 添加前缀（如“章节...”）
{}                           % 可用于在标题前插入代码
[{\color{CVBlue}\titlerule}]  % 在标题后插入一条横线
\titlespacing*{\section}{0cm}{*1.6}{*1.2}



\begin{document}
\pagenumbering{gobble}

%%%% 利用tikz来定位照片
\begin{tikzpicture}[remember picture, overlay] 
    \node[anchor = north east] at ($(current page.north east)+(-2cm,-1.2cm)$) {\includegraphics[height=3cm]{avatar.jpg}};
  \end{tikzpicture}%
  %%%% 利用tikz来定位学校Logo，这里只在第一页显示
  \begin{tikzpicture}[remember picture, overlay] 
    \node[anchor = north west] at ($(current page.north west)+(0.5cm,+1.0cm)$) {\includegraphics[height=6cm]{zju.png}};
  \end{tikzpicture}%
\centerline{\LARGE\bfseries{李映萱}} 

\centerline{\normalsize{\faPhone\ 111-1111-1111 \quad \faEnvelopeO\ \href{mailto:3230100000@zju.edu.cn}{3230100000@zju.edu.cn}}} 

\centerline{\normalsize{\faGithubSquare\ \href{https://github.com/maksymilan}{https://github.com/maksymilan} \quad \faRssSquare\ \href{https://maksymilan.github.io/}{https://maksymilan.github.io}}} 
    
\section{\makebox[\widthof{\faGraduationCap}][c]{\color{CVBlue}\faGraduationCap}\ 教育背景}    
\textbf{郑州大学} \hfill 2023.9 -- 至今\\[0.5em] % 标题和正文间加一点距离
人工智能\quad 大三
\begin{itemize}[nosep]
    \item 相关课程：概率论与数理统计（98）、线性代数（100）、离散数学（85）、高级语言程序设计（91）、深度学习（87）
\end{itemize}

\section{\makebox[\widthof{\faUsers}][c]{\color{CVBlue}\faUsers}\ 项目经历}

% --- 第一个项目 ---
% 将标题行末尾的 \\ 替换为 \titlebreak 命令
\textbf{SmartPPT智能会议图文助手} \hfill 2026.03 -- 2026.05 \titlebreak
项目描述： SmartPPT 是一款面向线下课堂、讲座和会议场景的智能图文笔记小程序，针对传统手动拍照、录音和整理笔记效率低、图文语音难关联的问题，构建了PPT自动采集、语音转写、OCR识别、AI摘要和PDF导出的完整流程。项目以PPT页面作为笔记结构核心，通过画面变化检测自动抓拍有效页面，并基于时间轴关联对应讲解音频，实现图文音同步整理和会后高效复盘。
% 在 itemize 的选项中，使用 topsep=\intraProjectListTopSep 来控制上边距
\begin{itemize}[nosep, topsep=\intraProjectListTopSep]
    \item \textbf{数据与后端}：后端采用Supabase构建BaaS架构，使用PostgreSQL设计profiles、folders、notes等数据表，并通过RLS策略实现用户数据隔离；使用Supabase Auth支持用户名、手机号及微信小程序登录，使用Supabase Storage存储图片和音频文件。AI 能力通过Supabase Edge Functions及第三方API集成实现，包括百度OCR、百度 ASR、文心大模型摘要和百度AI搜索，并在前端封装统一数据API，完成笔记CRUD、文件上传、AI调用和结构化Markdown内容生成。
    \item \textbf{可视化前端}：基于Taro 4.1、React 18和TypeScript开发微信小程序，使用Tailwind CSS与weapp-tailwindcss构建移动端界面。通过Taro Camera、RecorderManager、FileSystem、Canvas等微信小程序能力，实现PPT自动采集、语音录制与转写、识别结果编辑、笔记预览、音频播放和端侧PDF导出等功能。
    \item \textbf{核心产出}：全面掌握了从\textbf{数据采集、清洗、LLM标注到可视化分析}的“反烂坑”全链路技术。
\end{itemize}

% 使用 \projectsep 命令来分隔两个项目
\projectsep

% --- 第二个项目 ---
\textbf{AIGC驱动的超进化自动挖坑机 (Agent)} \hfill 2025.02 -- 至今 \titlebreak
项目描述：为探索心灵版规的边界，以及测试版主锁沉反应速度的极限，本项目基于\textbf{多智能体协作(MCP)与LLM}技术，开发了一款能\textbf{自主“挖坑”与“互坑”}的AI Agent。
\begin{itemize}[nosep, topsep=\intraProjectListTopSep]
    \item \textbf{并发与调度}：后端采用 \textbf{Go} 语言，充分利用 \textbf{Goroutine} 和 \textbf{Channel} 的高并发模型，模拟\textbf{用户并发在线}，执行定时挖坑，确保挖坑行动的\textbf{隐蔽性}与\textbf{高效性}。
    \item \textbf{实时监控与响应}：运用 \textbf{WebSocket} 协议实现了对目标帖子的\textbf{实时监听与交互}。确保在烂坑被识破的第一时间，Agent能光速完成\textbf{“lktp”}或\textbf{反向钓鱼}操作，展现出极高的AI博弈能力。
    \item \textbf{可视化GUI}：前端使用 \textbf{React} 框架构建了\textbf{后台管理的可视化面板 (Admin Panel)}。可一键下达挖坑指令、动态调整Agent的\textbf{指令}、并实时监控数据。
\end{itemize}

\section{\makebox[\widthof{\faCogs}][c]{\color{CVBlue}\faCogs}\ 技术栈}
\begin{itemize}[nosep]
    \item \textbf{编程语言：} \textbf{Go}, Python, C++
    \item \textbf{开发工具：} SSH, Git, Vim, MakeFile,LaTex
    \item \textbf{操作系统：} Linux
\end{itemize}
\section{\makebox[\widthof{\faGraduationCap}][c]{\color{CVBlue}\faList}\ 获奖情况}
\begin{itemize}
    \item 2025年全国大学生数学建模竞赛河南赛区省级一等奖 \hfill 2025.11
    \item 2025年全国大学生英语竞赛C类国家级三等奖 \hfill 2025.5
    \item 郑州大学2024-2025学年三好学生 \hfill 2025.11
    \item 郑州大学2023-2024学年优秀学生奖学金一等奖 \hfill 2024.12
\end{itemize}
    
\section{\makebox[\widthof{\faInfo}][c]{\color{CVBlue}\faInfo}\ 其他}
\begin{itemize}[parsep=0.5ex]
    \item \textbf{技术博客：} \href{https://maksymilan.github.io/}{https://maksymilan.github.io/}
    \item \textbf{GitHub：} \href{https://github.com/maksymilan}{https://github.com/maksymilan} 
    \item \textbf{英语水平：} CET-4, CET-6
\end{itemize}
\end{document}
